\documentclass[11pt]{article}
\usepackage{amsmath,amssymb,amsthm}
\usepackage[margin=1.2in]{geometry}
\newtheorem{theorem}{Theorem}
\newtheorem{lemma}[theorem]{Lemma}
\newtheorem{corollary}[theorem]{Corollary}
\newtheorem{proposition}[theorem]{Proposition}
\newtheorem{remark}[theorem]{Remark}
\title{Erd\H{o}s \#396: the antipodal law for shifted windows of
$\binom{2n}{n}$}
\author{Samuel Lavery}
\date{Draft --- \today}
\begin{document}
\maketitle

\begin{abstract}
\#396 asks whether for every $k$ there is an $n$ with
$\prod_{0\le i\le k}(n-i)\mid\binom{2n}{n}$.  We give an exact digit-sum criterion for
this divisibility, valid for every $n$ and $k$, and use it to prove an
\emph{antipodal law}: for every prime $p>\max(2k,\sqrt{2n})$ the residue $n\bmod p$ decides
both the divisor condition and the carry condition, and the two conditions occupy
complementary half-turns of the clock $\mathbb Z/p$.  Consequently a \emph{negative} shift
$n-i$ forces $p\nmid\binom{2n}{n}$, while a \emph{positive} shift $n+i$ forces
$p\mid\binom{2n}{n}$.  This proves that every solution of \#396 has all of
$n,n-1,\dots,n-k$ free of prime factors above $\max(2k,\sqrt{2n})$, and it explains why
Pomerance's companion problem for $\prod_{1\le i\le k}(n+i)$ has density $1$ while \#396
does not.  Two unconditional lower bounds on the least solution follow.
\#396 itself is not proved.
\end{abstract}

\section{Ledger form}

Throughout $p$ is prime, $s_p(m)$ is the sum of the base-$p$ digits of $m$, and
$v_p$ is the $p$-adic valuation.  We use Legendre's formula
$v_p(m!)=\bigl(m-s_p(m)\bigr)/(p-1)$ and Kummer's theorem in the form
\[
  v_p\binom{2n}{n}\;=\;\sum_{j\ge1}\Bigl(\bigl\lfloor 2n/p^{j}\bigr\rfloor
      -2\bigl\lfloor n/p^{j}\bigr\rfloor\Bigr)
  \;=\;\#\Bigl\{\,j\ge1:\ \bigl\{\,n/p^{j}\,\bigr\}\ \ge\ \tfrac12\,\Bigr\},
\]
each summand being $0$ or $1$.  We call the indicator $\{n/p^{j}\}\ge\frac12$ the
\emph{half-turn at level $j$}: it says that the point $n/p^{j}$ lies in the upper half of
the clock of circumference $1$, equivalently that $n\bmod p^{j}\ge p^{j}/2$.

\begin{theorem}[exact criterion]\label{thm:crit}
Let $n\ge k+1\ge1$.  Then
\[
  \prod_{0\le i\le k}(n-i)\ \Bigm|\ \binom{2n}{n}
  \qquad\Longleftrightarrow\qquad
  3\,s_p(n)\ \ge\ (k+1)+s_p(n-k-1)+s_p(2n)\ \ \text{for every prime }p .
\]
\end{theorem}

\begin{proof}
Since $\prod_{0\le i\le k}(n-i)=n!/(n-k-1)!$, Legendre's formula gives
\[
  v_p\!\left(\prod_{0\le i\le k}(n-i)\right)
  =\frac{\bigl(n-s_p(n)\bigr)-\bigl((n-k-1)-s_p(n-k-1)\bigr)}{p-1}
  =\frac{(k+1)-s_p(n)+s_p(n-k-1)}{p-1},
\]
and, again by Legendre,
\[
  v_p\binom{2n}{n}=\frac{\bigl(2n-s_p(2n)\bigr)-2\bigl(n-s_p(n)\bigr)}{p-1}
  =\frac{2s_p(n)-s_p(2n)}{p-1}.
\]
Divisibility holds if and only if the first quantity is at most the second for every $p$;
clearing the common positive denominator $p-1$ and rearranging gives the stated
inequality.
\end{proof}

\emph{Verified}: Theorem~\ref{thm:crit} agrees with direct evaluation of
$\binom{2n}{n}\bmod\prod_{0\le i\le k}(n-i)$ for all $1792$ pairs with $n<260$ and
$k\le6$ --- no mismatch.

The criterion is exact and involves no estimate: the whole of the divisibility is the
comparison of three digit sums against the window length $k+1$.  Note that at $p=2$ one has
$s_2(2n)=s_2(n)$ and the criterion collapses to a two-term inequality; this is exploited in
Corollary~\ref{cor:binary}.

\section{The antipodal law}

Fix $k$ and $n$ and put
\[
  y\ :=\ \max\bigl(2k,\ \sqrt{2n}\bigr).
\]
For a prime $p>y$ the entire ledger at $p$ is determined by the single residue
$n\bmod p$, because $p^{2}>2n$ kills every level $j\ge2$.  The following three statements
are then exact descriptions of three arcs of the clock $\mathbb Z/p$.

\begin{lemma}[the three arcs]\label{lem:arcs}
Let $p>y$ and write $\rho=n\bmod p$.  Then
\begin{enumerate}
\item[(a)] $p$ divides $\prod_{0\le i\le k}(n-i)$ if and only if $\rho\in[0,k]$, and then
      $v_p\bigl(\prod_{0\le i\le k}(n-i)\bigr)=1$;
\item[(b)] $p$ divides $\prod_{1\le i\le k}(n+i)$ if and only if $\rho\in[p-k,p-1]$, and
      then $v_p\bigl(\prod_{1\le i\le k}(n+i)\bigr)=1$;
\item[(c)] $p$ divides $\binom{2n}{n}$ if and only if $\rho\ge p/2$, and then
      $v_p\binom{2n}{n}=1$.
\end{enumerate}
\end{lemma}

\begin{proof}
(a) $p\mid n-i$ for some $0\le i\le k$ means $n\equiv i\pmod p$; as $0\le i\le k<p$ this is
exactly $\rho\in[0,k]$, and the $i$ is unique since $p>k$.  Moreover $p^{2}>2n>n-i$, so
$v_p(n-i)=1$.  (b) Likewise $p\mid n+i$ with $1\le i\le k$ means $\rho=p-i$, and
$p^{2}>2n\ge n+i$ gives $v_p(n+i)=1$.  (c) By Kummer, $v_p\binom{2n}{n}$ counts the levels
$j\ge1$ with $\{n/p^{j}\}\ge\frac12$.  For $j\ge2$ we have $p^{j}\ge p^{2}>2n>n$, so
$\{n/p^{j}\}=n/p^{j}<\frac12$ and the level contributes $0$.  Only $j=1$ survives, and
$\{n/p\}=\rho/p\ge\frac12$ is $\rho\ge p/2$.
\end{proof}

Since $p>2k$ we have $k<p/2$, so the arc $[0,k]$ of (a) and the arc $[p/2,p)$ of (c) are
\emph{disjoint}, while the arc $[p-k,p-1]$ of (b) is \emph{contained} in the arc of (c).
The two shift directions sit at antipodal phases of the same clock.  This is the whole
content of the next two theorems.

\begin{theorem}[negative shifts: extinction]\label{thm:neg}
Let $p>\max(2k,\sqrt{2n})$ be prime.  If $p$ divides $n-i$ for some $0\le i\le k$, then
$v_p\binom{2n}{n}=0$.  Consequently, if $\prod_{0\le i\le k}(n-i)\mid\binom{2n}{n}$ then
no prime exceeding $\max(2k,\sqrt{2n})$ divides any of $n,n-1,\dots,n-k$.
\end{theorem}

\begin{proof}
By Lemma~\ref{lem:arcs}(a) the hypothesis gives $\rho=n\bmod p\in[0,k]$, and $k<p/2$ because
$p>2k$; so $\rho<p/2$ and Lemma~\ref{lem:arcs}(c) gives $v_p\binom{2n}{n}=0$.  If the
divisibility held we would need $1=v_p(n-i)\le v_p\binom{2n}{n}=0$, a contradiction.
\end{proof}

\begin{corollary}[the smooth-window requirement]\label{cor:smooth}
If $\prod_{0\le i\le k}(n-i)\mid\binom{2n}{n}$ then every one of the $k+1$ consecutive
integers $n-k,\;n-k+1,\;\dots,\;n$ is $y$-smooth, where $y=\max\bigl(2k,\sqrt{2n}\bigr)$.
\end{corollary}

\begin{proof}
Immediate from Theorem~\ref{thm:neg}: a prime factor of any $n-i$ exceeding $y$ is
excluded.
\end{proof}

\emph{Verified}: over all solutions $(n,k)$ with $n<400$ and $k\le7$, all $111$ prime
factors of the $k+1$ window entries are $\le y$ --- no exception.  Separately, over
$n<1500$ and $k\le8$, all $42\,112$ prime factors exceeding $y$ of window entries $n-i$
occur at non-solutions and every one of them satisfies $v_p\binom{2n}{n}=0$, as
Theorem~\ref{thm:neg} requires.

\begin{theorem}[positive shifts: automatic]\label{thm:pos}
Let $p>\max(2k,\sqrt{2n})$ be prime and let $1\le i\le k$.  If $p$ divides $n+i$ then
$v_p(n+i)=1\le v_p\binom{2n}{n}$.  Consequently the divisibility
$\prod_{1\le i\le k}(n+i)\mid\binom{2n}{n}$ imposes no condition whatsoever at primes above
$\max(2k,\sqrt{2n})$: it is decided entirely by the primes $p\le\max(2k,\sqrt{2n})$, and in
particular it places no smoothness requirement on the window.
\end{theorem}

\begin{proof}
By Lemma~\ref{lem:arcs}(b), $\rho=n\bmod p=p-i$ and $v_p(n+i)=1$.  Since $i\le k<p/2$ we
get $\rho=p-i>p/2$, so Lemma~\ref{lem:arcs}(c) gives $v_p\binom{2n}{n}=1$.
\end{proof}

\emph{Verified}: over $n<1500$ and $k\le8$ there are $33\,752$ pairs $(n+i,p)$ with
$p>y$ and $p\mid n+i$; in every one $v_p(n+i)=1$ and $v_p\binom{2n}{n}=1$ --- no exception.

Theorems~\ref{thm:neg} and~\ref{thm:pos} are the same computation read at two antipodal
phases, and together they account for the asymmetry in the literature.  Pomerance showed
that the set of $n$ with $\prod_{1\le i\le k}(n+i)\mid\binom{2n}{n}$ has density $1$; by
Theorem~\ref{thm:pos} that statement is a statement about the finitely many primes below
$\max(2k,\sqrt{2n})$ only, where the digit-sum criterion of Theorem~\ref{thm:crit} is
satisfied for almost all $n$.  By contrast Corollary~\ref{cor:smooth} shows that \#396
cannot be attacked prime by prime at all: it demands a block of $k+1$ consecutive integers
with no prime factor above $\sqrt{2n}$, an event of density $\asymp\rho(2)^{k+1}$ in the
Dickman function, and this --- not any local carry condition --- is what makes \#396 hard.

\section{Two unconditional lower bounds on a solution}

\begin{corollary}[binary bound]\label{cor:binary}
If $\prod_{0\le i\le k}(n-i)\mid\binom{2n}{n}$ then
\[
  2\,s_2(n)\ \ge\ (k+1)+s_2(n-k-1),
  \qquad\text{hence}\qquad
  s_2(n)\ \ge\ \Bigl\lceil\tfrac{k+1}{2}\Bigr\rceil
  \quad\text{and}\quad n\ \ge\ 2^{(k-1)/2}.
\]
\end{corollary}

\begin{proof}
Apply Theorem~\ref{thm:crit} at $p=2$ and use $s_2(2n)=s_2(n)$ to cancel one copy of
$s_2(n)$.  The second inequality drops the nonnegative term $s_2(n-k-1)$, and the third
follows from $s_2(n)\le\lfloor\log_2 n\rfloor+1$.
\end{proof}

\emph{Verified}: all $329$ solutions with $n<3000$, $k\le7$ satisfy the two-term
inequality --- no exception.

\begin{corollary}[prime-count bound]\label{cor:picount}
If $\prod_{0\le i\le k}(n-i)\mid\binom{2n}{n}$ and $y=\max(2k,\sqrt{2n})$ then
\[
  \pi(y)\ \ge\ (k+1)\,\frac{\log(n-k)}{\log(2n)} .
\]
In particular $\pi(y)\ge k$ once $n\ge(2n)^{k/(k+1)}+k$, so $y$ is at least the $k$-th
prime; when the maximum defining $y$ is attained by $\sqrt{2n}$ this gives
$n\ge\frac12 p_k^{2}$.
\end{corollary}

\begin{proof}
By Corollary~\ref{cor:smooth} the product $\prod_{0\le i\le k}(n-i)$ is $y$-smooth, and it
divides $\binom{2n}{n}$, so it divides $\prod_{p\le y}p^{\,v_p\binom{2n}{n}}$.  Kummer's
theorem bounds each carry count by $v_p\binom{2n}{n}\le\lfloor\log_p(2n)\rfloor$, i.e.\
$p^{\,v_p\binom{2n}{n}}\le 2n$.  Hence
\[
  (n-k)^{k+1}\ <\ \prod_{0\le i\le k}(n-i)\ \le\ (2n)^{\pi(y)} ,
\]
and taking logarithms gives the claim.
\end{proof}

\emph{Verified}: at the least solution for $k=4$, namely $n=45153$, one has $y=300$,
$\pi(300)=62$ and the right-hand side is $4.70$.

\section{Census}

Let $n_k$ denote the least $n$ with $\prod_{0\le i\le k}(n-i)\mid\binom{2n}{n}$
(OEIS A375077).  Searching by the criterion of Theorem~\ref{thm:crit}, with the window
prefiltered by Corollary~\ref{cor:smooth}, gives
\[
  n_0=2,\quad n_1=210,\quad n_2=2480,\quad n_3=8178,\quad n_4=45153,
  \qquad n_5>3\times10^{6}.
\]
Every one of these has its whole window $y$-smooth, as Corollary~\ref{cor:smooth} forces:
for instance at $k=1$, $n_1=210=2\cdot3\cdot5\cdot7$ and $209=11\cdot19$, with
$y=\max(2,\lfloor\sqrt{420}\rfloor)=20$, and indeed $19\le20$.  The growth of $n_k$ is far
faster than the bounds of Corollaries~\ref{cor:binary} and~\ref{cor:picount}
($2^{(k-1)/2}$ is $2.8$ at $k=4$, against $45153$), which is consistent with
Corollary~\ref{cor:smooth}: the binding constraint is the smooth window, not any local
carry budget.

\begin{remark}[harmonic content]
The object's own scale is the clock $\mathbb Z/p^{j}$, never unit $1$; the DC mode is the
Legendre digit-sum identity, which is evaluated \emph{exactly} in Theorem~\ref{thm:crit}
with no remainder to estimate, and the entire AC content at a prime $p>y$ is the single
half-turn indicator $\mathbf 1[\,\rho\ge p/2\,]$.  The difficulty of \#396 lives on the
carrier --- the requirement that a whole window of length $k+1$ be $\sqrt{2n}$-smooth ---
and not in any chart readout.
\end{remark}

\begin{remark}[register]
\textbf{Proved}: Theorem~\ref{thm:crit} (exact digit-sum criterion, all $n,k$);
Lemma~\ref{lem:arcs}; Theorem~\ref{thm:neg} and Corollary~\ref{cor:smooth} (every solution
of \#396 has a $\max(2k,\sqrt{2n})$-smooth window); Theorem~\ref{thm:pos} (the positive-shift
problem has no large-prime condition at all); Corollaries~\ref{cor:binary}
and~\ref{cor:picount}.  All are elementary consequences of Legendre and Kummer, and all
were checked against direct computation before being written down.
\textbf{Measured}: $n_0,\dots,n_4$ and $n_5>3\times10^{6}$; the verification counts
recorded after each statement.
\textbf{Not proved}: \#396 itself.  What Corollary~\ref{cor:smooth} leaves is the
existence, for every $k$, of $k+1$ consecutive $\sqrt{2n}$-smooth integers additionally
satisfying the digit-sum criterion at every prime below $\sqrt{2n}$; no construction of
such blocks is given here, and none is claimed.
\textbf{Falsifier}: a solution $(n,k)$ with a prime factor of some $n-i$ exceeding
$\max(2k,\sqrt{2n})$ would refute Theorem~\ref{thm:neg}; none was found in
$42\,112$ tested large-prime incidences.
\end{remark}
\end{document}
